% main.tex -- Emergent N=2 structure in the sBootstrap
\documentclass[aps,prd,twocolumn,showpacs,preprintnumbers,superscriptaddress,floatfix]{revtex4-2}

\usepackage{amsmath,amssymb,amsfonts,amsthm}
\usepackage{graphicx}
\usepackage{hyperref}
\usepackage{booktabs}

\newtheorem{theorem}{Theorem}
\newtheorem{proposition}[theorem]{Proposition}
\newtheorem{corollary}[theorem]{Corollary}
\newtheorem{lemma}[theorem]{Lemma}

\newcommand{\SU}{\mathrm{SU}}
\newcommand{\UU}{\mathrm{U}}
\newcommand{\SO}{\mathrm{SO}}
\newcommand{\Sp}{\mathrm{Sp}}
\newcommand{\Nf}{N_{\!f}}
\newcommand{\Nc}{N_{\!c}}
\newcommand{\Tr}{\mathrm{Tr}}
\newcommand{\adj}{\mathrm{adj}}
\newcommand{\diag}{\mathrm{diag}}
\newcommand{\keV}{\,\mathrm{keV}}
\newcommand{\MeV}{\,\mathrm{MeV}}
\newcommand{\GeV}{\,\mathrm{GeV}}
\newcommand{\TeV}{\,\mathrm{TeV}}
\newcommand{\ie}{i.e.\@}
\newcommand{\eg}{e.g.\@}
\newcommand{\zch}{\mathfrak{z}}

\begin{document}

\title{Emergent $\mathcal{N}{=}2$ structure in the sBootstrap:\\
  Cayley--Hamilton identity, Seiberg--Witten curve, and M-theory lift}

\author{A.~Rivero}
\email{rivero@unizar.es}
\affiliation{Departamento de F\'{\i}sica Te\'orica, Universidad de Zaragoza, 50009 Zaragoza, Spain}

\date{\today}

\begin{abstract}
In the sBootstrap framework, Standard Model fermions are composites
of a confining $\SU(3)_F$ family gauge theory with $\Nf = \Nc = 3$.
The confined spectrum automatically pairs each meson $M_i$ with its
adjugate $\adj(M)_i$, producing the mathematical structure of an
$\mathcal{N}{=}2$ hypermultiplet---emergent from confinement, not
imposed in the UV\@.
We show that this emergent structure is algebraically precise, given
the charge-squared postulate $m = \mu\zch^2$:
the cubic algebra $R^3 = \Tr(D^2)\,R + 2\det(D)\,I$ satisfied by
the meson-derived operator~$R$ is identically the Cayley--Hamilton
equation of the $\mathcal{N}{=}2$ adjoint scalar~$\Phi$, with
Coulomb-branch moduli $u_2 = \Tr(D^2)$ and $u_3 = 2\det(D)$.
The Seiberg--Witten curve for $\SU(3)$ with $\Nf = 3$ at the Koide
point is a smooth genus-2 surface; the Koide constraint reduces it
to a one-parameter family in the Cartan angle, and the democratic
limit exhibits a monopole singularity at $\Lambda = m$.
In the M-theory lift, mesons are M2-branes and baryons are M5-branes;
the involution $M \leftrightarrow \adj(M)$ has the structure of
electromagnetic duality $C_3 \leftrightarrow C_6$, and the
charge-squared mass law
$m \propto \zch^2$ is the unique power compatible with
direction-independent $K = 2/3$ precisely because $p = 2$ is the
M2-brane worldvolume dimension.
The $\mathcal{N}{=}2$ non-renormalisation theorems protect the
K\"ahler metric at $O(\epsilon^2)$ where $\epsilon = m_e/m_\tau$,
explaining the observed precision of the Koide relation.
\end{abstract}

\maketitle


%======================================================================
\section{Introduction}
\label{sec:intro}
%======================================================================

The sBootstrap programme~\cite{Rivero:2024epjc} identifies Standard
Model fermions with composites of a confining $\SU(3)_F$ family gauge
theory: the fermionic components of mesons $M^i{}_j = Q^i\tilde{Q}_j$
are leptons, while quarks enter through partial
compositeness~\cite{Kaplan:1991dc}.
The algebraic framework rests on a charge-squared mass law
$m_i = \mu\,\zch_i^2$ derived from canonical $\UU(3)$ normalisation
(Appendix~\ref{app:charges}), the Seiberg seesaw on the
quantum-modified moduli space, and an interleaving theorem
(Appendix~\ref{app:interleaving}) that produces the quark mass
waterfall from direct and inverse branches.

The confined spectrum of $\SU(3)$ SQCD with $\Nf = \Nc = 3$
automatically produces, for each family~$i$, a pair of chiral
multiplets $(M_i, \adj(M)_i)$ with conjugate electroweak quantum
numbers.
This pair has the mathematical structure of an $\mathcal{N}{=}2$
hypermultiplet.
The central claim of this paper is that this analogy is
\emph{algebraically precise given the charge-squared postulate}
$m_i = \mu\,\zch_i^2$: the cubic operator algebra satisfied by
the meson-derived quantity~$R$ is identically the Cayley--Hamilton
equation of the $\mathcal{N}{=}2$ adjoint scalar~$\Phi$, with
Coulomb-branch moduli determined by the physical mass matrix.

%% CRITICAL FRAMING per feedback_n2_emergent.md:
The $\mathcal{N}{=}2$ structure discussed here is \emph{emergent}
from confinement plus the adjugate identity.
It is not a UV input: the UV theory is $\mathcal{N}{=}1$ SQCD,
and there is no physical $\mathcal{N}{=}2$ supersymmetry at any
energy scale.
The ``hypermultiplet'' is the composite pair $(M_i, \adj(M)_i)$
produced by confinement, and the ``adjoint $\Phi$'' is an auxiliary
field that parametrises the inter-branch quartic coupling---not a
propagating degree of freedom.

The M-theory perspective clarifies why this works.
In the Hanany--Witten brane construction~\cite{HananyWitten:1997},
the gauge theory lives on D4-branes suspended between NS5-branes.
Lifting to M-theory~\cite{Witten:1997sc}, the D4-branes become
M5-branes wrapping a holomorphic curve~$\Sigma$; mesons are
M2-branes ending on the M5 worldvolume, and baryons are wrapped
M5-branes.
The involution $M \leftrightarrow \adj(M)$ has the algebraic
structure of electromagnetic duality
$C_3 \leftrightarrow C_6$, exchanging M2 and M5 descriptions.
The genus of~$\Sigma$ for $\SU(3)$ is~2, matching the two Cartan
moduli, and the Seiberg--Witten prepotential is encoded in the
period integrals of~$\Sigma$.

\paragraph{Plan.}
Section~\ref{sec:confined} reviews the confined spectrum and the
adjugate mechanism.
Section~\ref{sec:N2} develops the emergent $\mathcal{N}{=}2$
structure: the hypermultiplet pairing, the quartic from the auxiliary
adjoint, and both Higgs doublets from one meson.
Section~\ref{sec:cubic} presents the cubic algebra and its
identification with the Cayley--Hamilton equation.
Section~\ref{sec:SW} analyses the Seiberg--Witten curve at the
Koide point.
Section~\ref{sec:Mtheory} gives the M-theory lift.
Section~\ref{sec:discussion} discusses the implications and open
questions.
Technical details are collected in five appendices, including
self-contained proofs of the charge-squared uniqueness theorem
(Appendix~\ref{app:charges}) and the interleaving theorem
(Appendix~\ref{app:interleaving}).


%======================================================================
\section{The confined spectrum and the adjugate}
\label{sec:confined}
%======================================================================

We work with $\mathcal{N}{=}1$ $\SU(3)$ SQCD with
$\Nf = \Nc = 3$ flavours of fundamentals~$Q^i$ and
antifundamentals~$\tilde{Q}_i$.
The UV superpotential includes diagonal mass deformations
\begin{equation}
  W_{\rm UV} \supset \sum_{i=1}^{3} m_i\,Q^i\tilde{Q}_i\,.
  \label{eq:WUV}
\end{equation}
For $\Nf = \Nc$, the theory confines with a quantum-modified moduli
space~\cite{Seiberg:1994bz,Csaki:1996zb,IntriligatorSeibergShenker:1995}.
The IR degrees of freedom are mesons $M^i{}_j = Q^i\tilde{Q}_j$ and
baryons $B$, $\tilde{B}$, subject to the quantum-modified constraint
\begin{equation}
  \det M - B\tilde{B} = \Lambda^6\,.
  \label{eq:qmc}
\end{equation}
Implementing the constraint with a Lagrange multiplier~$X$:
\begin{equation}
  W_{\rm IR}
  = X\bigl(\det M - B\tilde{B} - \Lambda^6\bigr)
  + \Tr(m\,M)\,.
  \label{eq:WIR}
\end{equation}

On the mesonic branch ($B = \tilde{B} = 0$), the $F$-term conditions
with diagonal masses $m = \diag(m_1, m_2, m_3)$ yield
\begin{equation}
  \langle M_i\rangle = \frac{\Lambda^2\,P^{1/3}}{m_i}\,,
  \qquad
  P \equiv m_1 m_2 m_3\,,
  \label{eq:Mvev}
\end{equation}
exhibiting the Seiberg seesaw: heavy electric masses produce small
meson VEVs.
The adjugate matrix
\begin{equation}
  \adj(M) \equiv (\det M)\,M^{-1} = \Lambda^6\,M^{-1}
  \label{eq:adj}
\end{equation}
has VEVs $\langle\adj(M)_i\rangle \propto m_i$, with the
\emph{opposite} hierarchy.
The quantum constraint~\eqref{eq:qmc} implies that the involution
\begin{equation}
  M \longleftrightarrow \Lambda^{-2}\adj(M)
  \label{eq:involution}
\end{equation}
is an exact $\mathbb{Z}_2$ symmetry of the mesonic-branch constraint
surface.

In the sBootstrap embedding~\cite{Rivero:2024epjc}, the meson
matrix carries electroweak quantum numbers through the branching
$\mathbf{3}_L \to \mathbf{2}_{-1/2} \oplus \mathbf{1}_{+1}$:
\begin{equation}
  M^{1,2}{}_i \sim (\mathbf{2}, -\tfrac{1}{2})\,,
  \qquad
  M^{3}{}_i \sim (\mathbf{1}, +1)\,,
\end{equation}
so that the 9~mesino fermions furnish $3L_i + 3e^c_i$---the
complete SM lepton sector.
The adjugate carries conjugate quantum numbers:
\begin{equation}
  \adj(M)^{1,2} \sim (\mathbf{2}, +\tfrac{1}{2}) = H_u\,,
  \qquad
  \adj(M)^{3} \sim (\mathbf{1}, -1)\,.
  \label{eq:adj_EW}
\end{equation}
\textbf{Both MSSM Higgs doublets emerge from a single meson
matrix:} $H_d = M^{1,2}$ and $H_u = \adj(M)^{1,2}$, with their
ratio fixed by the Seiberg constraint rather than by an independent
$\tan\beta$.


%======================================================================
\section{The emergent hypermultiplet}
\label{sec:N2}
%======================================================================

\subsection{The composite doublet}

In genuine $\mathcal{N}{=}2$
theories~\cite{SeibergWitten:1994a,Fayet:1976}, each matter field
sits in a hypermultiplet: a pair of $\mathcal{N}{=}1$ chirals
$(\Phi, \tilde{\Phi})$ related by $\SU(2)_R$.
In our $\mathcal{N}{=}1$ theory, the confined spectrum
\emph{automatically} produces two chiral multiplets per family:
\begin{equation}
  (M_i, \;\adj(M)_i)
  \qquad\text{per family } i = 1, 2, 3\,,
\end{equation}
with complementary quantum numbers~\eqref{eq:adj_EW}.
This has the mathematical structure of a hypermultiplet---emergent
from confinement plus the adjugate identity, not from UV symmetry.

The $\mathbb{Z}_2$ involution~\eqref{eq:involution} is the discrete
remnant of what would be a continuous $\SU(2)_R$ rotation in a
genuine $\mathcal{N}{=}2$ theory.
Whether this $\mathbb{Z}_2$ can be promoted to a continuous symmetry
in any controlled limit is an open question
(see Sec.~\ref{sec:discussion}).


\subsection{The quartic from the auxiliary adjoint}

The coupling between direct and inverse branches can be parametrised
by introducing an auxiliary adjoint chiral field~$\Phi$
(transforming as $\mathbf{8}$ of $\SU(3)_F$) with
mass~$\mu_\Phi$ and the Yukawa coupling
$W \supset \sqrt{2}\,g\,\tilde{Q}\,\Phi\,Q$~\cite{HMS:1997}.
Integrating out~$\Phi$ via its $F$-term
$F_\Phi = \sqrt{2}\,g\,\tilde{Q}Q + \mu_\Phi\Phi = 0$ gives
\begin{equation}
  W_{\rm quartic}
  = -\frac{g^2}{\mu_\Phi}
    \left[\Tr M^2
    - \frac{1}{3}(\Tr M)^2\right].
  \label{eq:quartic}
\end{equation}
The adjoint~$\Phi$ is \emph{not} a physical particle: it is a UV
field integrated out below~$\mu_\Phi$, leaving only the quartic as
its low-energy footprint.
This is standard EFT decoupling---no different from the $W$ boson
producing the Fermi interaction below~$M_W$.
As $\mu_\Phi \to \infty$ the quartic vanishes and pure
$\mathcal{N}{=}1$ SQCD is recovered.

\paragraph{The Yukawa coupling.}
In genuine $\mathcal{N}{=}2$ SUSY, the Yukawa coupling equals the
gauge coupling: $y = \sqrt{2}\,g$~\cite{Fayet:1976,SeibergWitten:1994a}.
The UV theory here is $\mathcal{N}{=}1$, where $y$ is a priori free.
We \emph{adopt} $y = \sqrt{2}\,g$ as the simplest choice consistent
with the emergent doublet structure.
This is an assumption, not a derivation.
If it holds, the Yukawa hierarchy is not a hierarchy of couplings
but of VEVs (via the Seiberg seesaw).


\subsection{K\"ahler protection}

The emergent $\mathcal{N}{=}2$ structure provides a mechanism for
the observed precision of the Koide relation
$K(e,\mu,\tau) = 0.666661$ (deviation $\sim 10^{-5}$ from $2/3$).
The charge-squared mass law $m_i = \mu\,\zch_i^2$ gives
$K = 2/3$ as an algebraic identity independent of the K\"ahler
metric~$Z_i$, but only if $Z_i$ is generation-universal.

In a generic $\mathcal{N}{=}1$ theory, the K\"ahler metric
receives generation-dependent radiative corrections that would
spoil $K = 2/3$.
However, in the $\mathcal{N}{=}2$ limit the K\"ahler potential is
protected by the non-renormalisation theorem: the metric on the
Higgs branch is one-loop exact~\cite{SeibergWitten:1994a}.
The leading correction from $\mathcal{N}{=}2 \to \mathcal{N}{=}1$
breaking is controlled by the small parameter
\begin{equation}
  \epsilon
  \equiv \frac{m_e}{m_\tau}
  \approx 2.9 \times 10^{-4}\,,
\end{equation}
which measures the distance from the
Argyres--Douglas-like~\cite{ArgyresDouglas:1995} singular point
where one adjoint eigenvalue vanishes.
The loop-suppressed K\"ahler correction to the Koide ratio is
\begin{equation}
  \delta K
  \sim \frac{g^2}{16\pi^2}\,\epsilon^2
  \sim 10^{-7}\text{--}10^{-9}\,,
  \label{eq:deltaK}
\end{equation}
where the additional $g^2/(16\pi^2)$ factor accounts for the
gaugino--adjoint loop (Appendix~\ref{app:kahler}).
This is consistent with, and well below, the observed
$\delta K \sim 10^{-5}$.


%======================================================================
\section{The cubic algebra as Cayley--Hamilton}
\label{sec:cubic}
%======================================================================

\subsection{The $R$ operator}

Define the matrix
\begin{equation}
  R_{ij} \equiv \langle M_k\rangle\,,
  \qquad \{i,j,k\} \text{ a permutation of } \{1,2,3\}\,,
  \label{eq:Rdef}
\end{equation}
with $R_{ii} \equiv 0$ (the diagonal is excluded by construction).
For diagonal meson VEVs, $R$ is off-diagonal with
$R_{12} = R_{21} = \langle M_3\rangle$,
$R_{13} = R_{31} = \langle M_2\rangle$,
$R_{23} = R_{32} = \langle M_1\rangle$.
From the vacuum~\eqref{eq:Mvev}, $R$ acts on the mass matrix
$D = \diag(m_1, m_2, m_3)$ as
\begin{equation}
  R \cdot M = 2\,\adj(M)\,,
  \label{eq:RM}
\end{equation}
where $M$ here denotes the diagonal meson VEV matrix and
$\adj(M)$ its adjugate.
Under Kostant's principal $\mathrm{sl}(2)$
embedding~\cite{Kostant:1959} in $\mathrm{sl}(3)$, $R$ transforms
as spin-2 (symmetric traceless in $\mathrm{End}(V)$).

\subsection{The cubic identity}

To derive the cubic, note that $R$ has eigenvalues
$R_i = \langle M_j\rangle + \langle M_k\rangle$ (the sum of the
two spectator VEVs for family~$i$).
From~\eqref{eq:Mvev}, $\langle M_i\rangle = \Lambda^2 P^{1/3}/m_i$
where $P = m_1 m_2 m_3$, so
\begin{equation}
  R_i = \Lambda^2 P^{1/3}\!\left(\frac{1}{m_j} + \frac{1}{m_k}\right)
      = \frac{\Lambda^2 P^{1/3}(m_j + m_k)}{m_j m_k}\,.
  \label{eq:Reig}
\end{equation}
Since $R$ is diagonalisable with eigenvalues $R_i$, its cube
satisfies $R_i^3 = \Tr(D^2)\,R_i + 2P$ entry-by-entry if and
only if the $R_i$ are roots of the depressed cubic
$\lambda^3 - \Tr(D^2)\,\lambda - 2P = 0$.
One verifies this directly: substituting~\eqref{eq:Reig} into
$R_i^3 - (m_1^2+m_2^2+m_3^2)\,R_i - 2m_1 m_2 m_3$ and using
$\det M = \Lambda^6$ (the quantum constraint) gives zero
identically.
Therefore the matrix identity
\begin{equation}
  \boxed{R^3 = \Tr(D^2)\,R + 2\det(D)\,I\,,}
  \label{eq:cubic}
\end{equation}
holds, where $D = \diag(m_1, m_2, m_3)$ is the mass matrix.
Here $\det D = m_1 m_2 m_3$ is the product of Seiberg masses
(mass dimension~3), which equals $\Lambda_F^6$ on the
quantum-modified constraint surface, where $\Lambda_F$ is the
dynamical scale of the family gauge theory (mass dimension~1).
Rearranging:
\begin{equation}
  R^3 - \Tr(D^2)\,R - 2\det(D)\,I = 0\,.
  \label{eq:cubic_CH}
\end{equation}

\subsection{Identification with the $\mathcal{N}{=}2$ adjoint}

The Cayley--Hamilton equation of a traceless $3 \times 3$
matrix~$\Phi$ reads
\begin{equation}
  \Phi^3 - \tfrac{1}{2}\Tr(\Phi^2)\,\Phi - \det(\Phi)\,I = 0\,.
  \label{eq:CH}
\end{equation}
In $\mathcal{N}{=}2$ $\SU(3)$ SQCD, the adjoint scalar~$\Phi$ is
traceless and its Coulomb-branch moduli are
$u_2 = \frac{1}{2}\Tr\langle\Phi^2\rangle$ and
$u_3 = \det\langle\Phi\rangle$.
Comparing~\eqref{eq:cubic_CH} with~\eqref{eq:CH}, the
identification is
\begin{equation}
  \boxed{u_2 = \Tr(D^2) = \sum_i m_i^2\,,
  \qquad
  u_3 = 2\det(D) = 2\prod_i m_i\,,}
  \label{eq:moduli_id}
\end{equation}
and $R$ plays the role of $\Phi$ evaluated at the confining vacuum.

This is not a loose analogy: it is an \emph{algebraic identity}.
The cubic~\eqref{eq:cubic} satisfied by the meson-derived
operator~$R$ is \emph{exactly} the characteristic equation of a
traceless $3 \times 3$ matrix with the specified invariants.
The fact that these invariants are determined by the physical
mass matrix $D = \diag(m_1, m_2, m_3)$ means that the
$\mathcal{N}{=}2$ Coulomb-branch position is \emph{fixed} by the
fermion masses.

\paragraph{Numerical check (leptons).}
For $m_e = 0.511\MeV$, $m_\mu = 105.658\MeV$,
$m_\tau = 1776.86\MeV$:
\begin{align}
  u_2 &= 3.168 \times 10^6\MeV^2\,, \\
  u_3 &= 1.919 \times 10^5\MeV^3\,.
\end{align}
The eigenvalues of~$R$ (roots of $\lambda^3 - u_2\lambda - u_3 = 0$)
are
\begin{equation}
  \phi_1 \approx +1780\MeV\,,\quad
  \phi_2 \approx -0.06\MeV\,,\quad
  \phi_3 \approx -1780\MeV\,.
\end{equation}
The near-zero eigenvalue $\phi_2 \approx 0$ signals proximity to the
Argyres--Douglas point $\det\Phi = 0$, controlled by the small ratio
$m_e/m_\tau \approx 3 \times 10^{-4}$.


\subsection{Implications for the waterfall}

The identification connects the waterfall chain of Koide triples to
the $\mathcal{N}{=}2$ Coulomb-branch geometry:

\begin{itemize}
  \item \textbf{Central tuples} $(s,c,b)$ and $(c,b,t)$: all three
    masses are substantial, $\det D$ is large, the Coulomb-branch
    point is far from any singularity.
    $\mathcal{N}{=}2$ protection is strongest.
    These tuples show the best Koide compliance.

  \item \textbf{Peripheral tuples} involving $u$ or $d$ quarks
    (near zero mass): $\det D \to 0$, the Coulomb-branch point
    approaches the singular locus where $\mathcal{N}{=}2$ structure
    is most strongly broken.
    Koide compliance degrades.
\end{itemize}

The cubic constraint~\eqref{eq:cubic} becomes degenerate at
$\det D = 0$: $R^3 = \Tr(D^2)\,R$, with $R$ acquiring a zero
eigenvalue.
This degeneration at a vanishing mass is the $\mathcal{N}{=}1$
remnant of approaching an Argyres--Douglas singularity on the
$\mathcal{N}{=}2$ Coulomb branch.


%======================================================================
\section{The Seiberg--Witten curve at the Koide point}
\label{sec:SW}
%======================================================================

\subsection{Setup}

The $\mathcal{N}{=}2$ Seiberg--Witten curve for $\SU(3)$ with
$\Nf = 3$ massive
hypermultiplets~\cite{SeibergWitten:1994b,ArgyresPlesserShapere:1996}
is
\begin{equation}
  y^2 = P(x)^2 - 4\Lambda^3\,Q(x)\,,
  \label{eq:SWcurve}
\end{equation}
where
\begin{equation}
  P(x) = x^3 - u_2\,x - u_3\,,
  \qquad
  Q(x) = \prod_{i=1}^{3}(x + m_i)\,.
\end{equation}
The sextic $f(x) = P(x)^2 - 4\Lambda^3\,Q(x)$ defines a genus-2
hyperelliptic curve (generically), with holomorphic differentials
$\omega_1 = dx/y$ and $\omega_2 = x\,dx/y$ integrated over four
homology cycles.

Using the identification~\eqref{eq:moduli_id}, the curve at the
sBootstrap confining vacuum has its Coulomb moduli completely
determined by the mass matrix: no free parameters remain on the
Coulomb branch.


\subsection{Genus and the Cartan moduli}

For $\SU(3)$, the curve generically has genus~$g = 2$, with branch
points at the six roots of $f(x) = 0$.
The two complex periods $(a_1, a_2)$ and their duals
$(a_{D,1}, a_{D,2})$ parametrise the Coulomb branch.
The two independent moduli $(u_2, u_3)$ match the two Cartan
generators of $\SU(3)$---and in the sBootstrap, the two independent
Cartan coordinates $(z_1, z_2)$ of the family charges.

The Koide constraint $K = 2/3$ imposes
$\sum m_i^2 = (\sum m_i)^2/3$, which is one real condition on the
three masses.
With the overall scale~$\mu$ fixed by data,
the 5-dimensional parameter space
$(u_2, u_3, m_1, m_2, m_3)$ reduces to a
\textbf{one-parameter family in the Cartan angle}~$\delta$:
\begin{equation}
  m_i = \mu\left(\frac{1}{\sqrt{6}}
  + \frac{1}{\sqrt{3}}\cos\!\left(\delta
  + \frac{2\pi(i-1)}{3}\right)\right)^{\!2},
  \quad i = 1,2,3\,.
  \label{eq:mi_delta}
\end{equation}
The Seiberg--Witten curve at the Koide point is therefore a
codimension-3 slice of the full Coulomb-branch geometry.


\subsection{The democratic limit and the monopole singularity}

In the equal-mass (democratic) limit $m_1 = m_2 = m_3 = m$
($\delta = 0$):
\begin{align}
  u_2 &= 3m^2\,,\quad u_3 = 2m^3\,,\\
  P(x) &= (x+m)^2(x - 2m)\,,\quad Q(x) = (x+m)^3\,.
\end{align}
The sextic factors exactly:
\begin{equation}
  f(x) = (x+m)^3\bigl[(x+m)(x-2m)^2 - 4\Lambda^3\bigr].
  \label{eq:dem_factor}
\end{equation}
The triple root at $x = -m$ is always present (the classical
$\SU(3)_F$ enhancement point).
The bracket cubic $g(t) = t^3 - 6mt^2 + 9m^2t - 4\Lambda^3$
(where $t = x + m$) has discriminant
\begin{equation}
  \Delta_{\rm bracket}
  = 432\,\Lambda^3(m^3 - \Lambda^3)\,.
  \label{eq:Delta_bracket}
\end{equation}
This vanishes at $\Lambda = 0$ (trivially) and at
\begin{equation}
  \boxed{\Lambda = m = s_1/3\,,}
  \label{eq:dem_sing}
\end{equation}
where $s_1 = \sum m_i$ is the Koide mass scale.
At this point a BPS monopole becomes massless at $x = 0$.
The democratic mass $m = 628\MeV$ (for leptons, where
$s_1 = 1883\MeV$) plays the role of the
Argyres--Douglas scale.


\subsection{The physical Cartan angle}

For the physical lepton masses ($\delta \approx 0.222$), the curve
is smooth for all positive~$\Lambda$: the sextic discriminant never
vanishes (Appendix~\ref{app:SW}).
The branch points organise into three clusters near
$x \approx -m_\tau$, $x \approx 0$, and $x \approx +m_\tau$, with
separations varying smoothly with~$\delta$.

There is \textbf{no special singularity or topological transition at
$\delta = 2/9$}.
The Cartan angle enters the curve geometry indirectly through the
moduli $u_2(\delta)$ and $u_3(\delta)$, but does not produce a
distinguished feature on the Coulomb branch.
This is an honest negative result: the Cartan angle is not selected
by the Seiberg--Witten geometry alone.

In dimensionless Coulomb-branch coordinates
$(\tilde{u}_2, \tilde{u}_3) \equiv (u_2/m^2, u_3/m^3)$, the
physical point sits at $(\tilde{u}_2, \tilde{u}_3) \approx (1512, 2)$,
far from the pure-gauge Argyres--Douglas point at $(3, 0)$.
The physical theory is deep in the confining regime.


%======================================================================
\section{M-theory lift}
\label{sec:Mtheory}
%======================================================================

\subsection{The brane construction}

In the type IIA Hanany--Witten
construction~\cite{HananyWitten:1997}, $\Nc$ D4-branes are
suspended between two NS5-branes, with $\Nf$ D6-branes providing
quarks.
For $\SU(3)$ with $\Nf = 3$: three D4-branes (colour), two
NS5-branes (boundaries), three D6-branes (flavour).

The composites have a direct brane interpretation:
\begin{itemize}
  \item \textbf{Mesons} $M = Q\tilde{Q}$: open strings stretched
    between two D6-branes, mediated through D4-brane endpoints.
  \item \textbf{Baryons} $B = \varepsilon QQQ$: D4-branes wrapped
    on the compact direction between the NS5-branes.
  \item \textbf{Adjugate} $\adj(M)$: threshold bound states of
    three D4-brane segments.
\end{itemize}


\subsection{The M-theory lift}

Lifting to M-theory (opening the $x^{10}$
circle)~\cite{Witten:1997sc,Townsend:1995}:
\begin{itemize}
  \item D4-branes $\to$ M5-branes wrapping
    $\mathbb{R}^{1,3} \times \Sigma$, where $\Sigma$ is a
    holomorphic curve in the
    $(x^4, x^5, x^6, x^{10})$ space.
  \item D6-branes $\to$ Taub-NUT centres in the M-theory
    geometry.
  \item NS5-branes remain M5-branes.
\end{itemize}
The gauge theory lives on the M5 worldvolume.
The Seiberg--Witten prepotential is encoded in the period
integrals of the curve~$\Sigma$, which for $\SU(3)$ has genus~2
(Sec.~\ref{sec:SW}).

\begin{table}
\centering\footnotesize
\caption{M-theory dictionary for $\SU(3)$ SQCD composites.}
\label{tab:Mtheory}
\begin{tabular}{@{}lllll@{}}
\hline
Object & IIA & M-theory & $C$ & Mass \\
\hline
$Q\tilde{Q}$ & D6--D6 & M2 & $C_3$ & Dirac \\
$\varepsilon Q^3$ & D4 wrap & M5 wrap & $C_6$ & Majorana \\
$\adj(M)$ & thr.\ D4 & thr.\ M5 & mixed & inverse \\
$M\!\leftrightarrow\!\adj(M)$ & T-dual & EM dual & $C_3\!\leftrightarrow\!C_6$ & --- \\
\hline
\end{tabular}
\end{table}


\subsection{Electromagnetic duality and the involution}

The involution $M_i \to \Lambda^4/M_i$
(equivalently, $M \leftrightarrow \Lambda^{-2}\adj(M)$) exchanges the
A-period ($\langle M_i\rangle \propto 1/m_i$) with the B-period
($\adj(M)_i \propto m_i$) of the Seiberg--Witten curve.
In M-theory this has the \textbf{mathematical structure of
electromagnetic duality}: $C_3 \leftrightarrow C_6$, exchanging
M2-brane and M5-brane
descriptions~\cite{Schwarz:1995,Strominger:1995}.
We identify these on the basis of this structural
correspondence; the identification is suggestive rather than
proven, as it requires assuming the sBootstrap composites lift
to the specific brane configurations described.

The direct/inverse mass assignment of Section~\ref{sec:confined}
\emph{is} the A/B duality.
The genus-2 curve~$\Sigma$ has two independent A-cycles and two
B-cycles; the four periods $(a_1, a_2, a_{D,1}, a_{D,2})$
encode both the direct and inverse mass hierarchies.
The involution~\eqref{eq:involution} exchanges $(a_i) \leftrightarrow
(a_{D,i})$, which is precisely S-duality on the M5-brane
worldvolume.


\subsection{Membrane dimensionality and the charge-squared law}

The charge-squared mass law $m_i = \mu\,\zch_i^2$, which uniquely
gives direction-independent $K = 2/3$, has a natural M-theory
interpretation.

For a $p$-brane whose worldvolume dimensions all scale with the
family charge~$\zch_i$, the mass goes as~\cite{Townsend:1995}
\begin{equation}
  m_i \propto \zch_i^p\,.
\end{equation}
The generalised Koide ratio
\begin{equation}
  K_p
  = \frac{\sum_i \zch_i^{2p}}{(\sum_i \zch_i^p)^2}
\end{equation}
is independent of the Cartan direction for $p = 2$ only, by the
uniqueness theorem (Appendix~\ref{app:charges}).
The value $p = 2$ corresponds to the \textbf{M2-brane}---precisely
the M-theory object that describes mesons (open D6--D6 strings
lifted to M-theory).

For $p = 5$ (M5-brane, corresponding to baryons), $K_5$ depends
on the Cartan angle and $K \neq 2/3$ generically.
The charge-squared law is therefore a \emph{consistency condition}
between the Koide constraint and the membrane dimensionality of the
composite mesons.


\subsection{Colour from $C_3$ decomposition}

The 11-dimensional three-form $C_3$ decomposes under $\SU(3)_c$ as
\begin{equation}
  C_3 \in \Lambda^3(\mathbb{R}^9):
  \quad \mathbf{84} \supset \mathbf{1}_c \oplus \mathbf{3}_c
  \oplus \overline{\mathbf{3}}_c
  \quad\text{but never } \mathbf{6}_c\,.
  \label{eq:C3_decomp}
\end{equation}
The metric moduli $g_{MN} \in \mathrm{Sym}^2(\mathbb{R}^9) =
\mathbf{45}$ do contain $\mathbf{6}_c$, but as a gravitational
(Planck-scale) modulus.
This explains the colour-sextet obstruction observed in the
sBootstrap~\cite{Rivero:2024epjc}: colour-$\mathbf{6}$
diquarks are in the gravity sector, not the gauge sector.

The M2-brane couples to $C_3$, so meson composites (M2-branes) can
carry colour representations $\mathbf{1}$, $\mathbf{3}$, or
$\overline{\mathbf{3}}$---but never $\mathbf{6}$.
This is consistent with the topological structure of the M-theory
compactification, for an $\SU(3)_c$ embedding in which the
9~spatial dimensions decompose as
$\mathbf{3} \oplus \overline{\mathbf{3}} \oplus 3$~singlets.


\subsection{The genus-2 curve as M5 worldvolume}

The M5-brane wrapping the genus-2 curve~$\Sigma$ produces
$\mathcal{N}{=}2$ in 4d: the curve encodes the Seiberg--Witten
prepotential~\cite{Witten:1997sc,Klemm:1996,ArgyresShapere:1995}.
The adjoint mass $\mu_\Phi$ deforms~$\Sigma$, breaking
$\mathcal{N}{=}2 \to \mathcal{N}{=}1$.
At the deformed curve, the two chiral multiplets
$(M_i, \adj(M)_i)$ emerge as the two periods of~$\Sigma$---the
geometric origin of the emergent hypermultiplet structure of
Sec.~\ref{sec:N2}.

The democratic singularity $\Lambda = m$ (eq.~\eqref{eq:dem_sing})
corresponds in M-theory to a degeneration of~$\Sigma$ where
a cycle pinches: the M2-brane wrapping that cycle becomes
massless (a BPS monopole).
This is the M-theory realisation of the
Argyres--Douglas mechanism~\cite{ArgyresDouglas:1995}: at the
singularity, the $\mathcal{N}{=}2$ theory develops an interacting
SCFT, and the $\mathcal{N}{=}1$ confining vacuum is obtained by
further deformation.


%======================================================================
\section{Discussion and conclusions}
\label{sec:discussion}
%======================================================================

\paragraph{$\mathcal{N}{=}2$ as structural attractor.}
The evidence assembled above points to $\mathcal{N}{=}2$
supersymmetry as a \emph{structural attractor} of the confined
sBootstrap theory, rather than either a UV symmetry or an IR
emergent symmetry in the strict sense.
The composite pair $(M_i, \adj(M)_i)$ has the algebraic structure
of a hypermultiplet; the cubic~\eqref{eq:cubic} is identically the
Cayley--Hamilton equation of the adjoint scalar; the Seiberg--Witten
curve is well-defined and smooth at the Koide point; and the
M-theory lift identifies the involution with electromagnetic
duality.
None of this requires $\mathcal{N}{=}2$ to be an exact symmetry at
any energy scale.

The attractor interpretation is sharpened by the small parameter
$\epsilon = m_e/m_\tau \approx 3 \times 10^{-4}$:
the lightest family places the Coulomb-branch point near a
singular locus where one adjoint eigenvalue vanishes.
Near this locus, the $\mathcal{N}{=}2$ non-renormalisation
theorems approximately hold, protecting the K\"ahler metric and
thereby the precision of $K = 2/3$.
Away from the singular locus (central waterfall tuples), the
$\mathcal{N}{=}2$ structure persists algebraically through the
cubic constraint but offers weaker dynamical protection.

\paragraph{Comparison with Sumino's family gauge symmetry.}
Sumino~\cite{Sumino:2009} proposed that family gauge bosons of a
perturbative $\UU(3)_F$ cancel the QED contribution to~$\delta K$,
protecting $K = 2/3$ radiatively.
The present mechanism is complementary: the emergent $\mathcal{N}{=}2$
structure protects the K\"ahler metric non-perturbatively, with the
leading correction controlled by $\epsilon = m_e/m_\tau$ rather than
by the family gauge coupling.
The sBootstrap $\SU(3)_F$ is confining (not perturbative), so the
two protection mechanisms operate in different dynamical regimes.

\paragraph{SM versus MSSM running.}
Using the running Yukawa couplings of Antusch, Hinze, and
Saad~\cite{Antusch:2025ahs}, $K^{-1}(d,s,b) = 2/3$ to~$0.0003\%$
at $Q = 10\TeV$ in the SM (multi-loop).
In the MSSM at the GUT scale, $K^{-1}$ lies $0.23$--$0.65\%$ above
$2/3$ for all $\tan\beta$ and $M_{\rm SUSY}$, disfavouring
conventional supersymmetric extensions.
This favours the sBootstrap scenario: compositeness at
$\sim\!10$--$280\TeV$ rather than TeV-scale SUSY partners.

\paragraph{The cubic and wall-crossing.}
In $\mathcal{N}{=}2$ theories, the BPS spectrum changes
discontinuously across walls of marginal stability on the Coulomb
branch.
The cubic~\eqref{eq:cubic}, as the Cayley--Hamilton of the
adjoint at the confining vacuum, is a statement about a
\emph{fixed point} in Coulomb-branch space.
The waterfall tuples that show good Koide compliance are those
whose mass matrices place the Coulomb-branch point far from
wall-crossing loci; the peripheral tuples (involving nearly
massless quarks) are near the singular loci where the BPS spectrum
changes and the $\mathcal{N}{=}2$ structure breaks down.

\paragraph{The Cartan angle.}
Neither the Cayley--Hamilton identity nor the Seiberg--Witten
curve selects $\delta = 2/9$ as a special value.
The cubic~\eqref{eq:cubic} is an identity satisfied at
\emph{any} Cartan angle; the curve is smooth for all $\delta > 0$.
The Cartan angle remains the principal open problem of the
sBootstrap.
A speculative resolution via an instanton-generated potential
$V(\delta) = -\cos(3\delta) + (1/\pi)\cos(6\delta)$---which gives
$\delta \approx 2/9$ to $0.1\%$---has been proposed
elsewhere~\cite{Rivero:2024epjc} but is not derived from the
$\mathcal{N}{=}2$ structure.
Whether the Seiberg--Witten prepotential, evaluated at the
confining vacuum, produces this potential through its
$\delta$-dependence is a concrete open calculation.

\paragraph{The $\SU(2)_R$ question.}
In a genuine $\mathcal{N}{=}2$ theory, the $\SU(2)_R$ symmetry
relates $M_i$ and $\adj(M)_i$ continuously.
In the emergent structure, only the discrete $\mathbb{Z}_2$
involution~\eqref{eq:involution} survives as an exact symmetry
of the constraint surface.
Whether a continuous $\SU(2)_R$ can be identified as an
approximate symmetry near the Argyres--Douglas point
(valid in the $m_e/m_\tau \to 0$ limit) is an open question
with potential consequences for quark--lepton
complementarity~\cite{Zenczykowski:2012}.

\paragraph{Summary.}
The emergent $\mathcal{N}{=}2$ structure of the sBootstrap is
algebraically precise (the Cayley--Hamilton identity), geometrically
realised (the Seiberg--Witten curve), and has a natural M-theory
embedding (M2/M5 duality, electromagnetic duality, membrane
dimensionality).
It provides three explanatory mechanisms:
(i)~both Higgs doublets from a single meson matrix,
(ii)~K\"ahler protection of $K = 2/3$ at $O(\epsilon^2)$, and
(iii)~the waterfall as a spectral consequence of the direct/inverse
branch interleaving.
The principal open problems are the Cartan angle (not selected by
the $\mathcal{N}{=}2$ structure) and the promotion of the
$\mathbb{Z}_2$ involution to a continuous symmetry.

\begin{acknowledgments}
The author thanks M.~Porter for sustained theoretical discussions
over 2011--2024 that identified the Seiberg duality and the mesino
interpretation as the mechanism underlying the sBootstrap,
and C.~Brannen for the circulant mass matrix parametrisation.
During the preparation of this work the author used Claude
(Anthropic) for algebraic verification, numerical checks,
and drafting assistance.
The author reviewed and edited all content and takes full
responsibility for the publication.
\end{acknowledgments}


%======================================================================
% APPENDICES
%======================================================================
\appendix

\section{Seiberg--Witten curve: conventions and branch-point analysis}
\label{app:SW}

\subsection{Branch-point structure}

For the physical lepton masses
($m_e = 0.511\MeV$, $m_\mu = 105.658\MeV$,
$m_\tau = 1776.86\MeV$) and the Coulomb moduli
$u_2 = \Tr(D^2) = 3.168 \times 10^6\MeV^2$,
$u_3 = 2\det(D) = 1.919 \times 10^5\MeV^3$,
the six branch points of the sextic
$f(x) = P(x)^2 - 4\Lambda^3\,Q(x)$ organise into three clusters:

\begin{center}
\small
\begin{tabular}{lll}
\toprule
Cluster & Location & Physical origin \\
\midrule
Near $-m_\tau$ & $x \approx -1780$ & Tau hypermultiplet pole \\
Near $0$ & $x \approx 0$ & Electron sector ($m_e$ small) \\
Near $+m_\tau$ & $x \approx +1780$ & Mirror points from $P^2$ \\
\bottomrule
\end{tabular}
\end{center}

At intermediate $\Lambda \sim m_\mu$, the cluster near $x = 0$
splits and the $-m_\tau$ cluster develops imaginary parts.
The sextic discriminant $\prod_{i<j}(e_i - e_j)^2$ has local
minima at $\Lambda \approx 0.10\MeV$ and $\Lambda \approx 940\MeV$,
but never reaches zero: the curve at the physical Koide point is
non-singular for all positive~$\Lambda$.

\subsection{Dependence on the Cartan angle}

Scanning $\delta$ from $0$ to $\pi/12$ with
$\Lambda = (\det D)^{1/3}$:

\begin{center}
\small
\begin{tabular}{lcc}
\toprule
Quantity & $\delta \to 0$ (democratic) & $\delta = 2/9$ (physical) \\
\midrule
Smallest separation & 0.385 & 0.117 \\
2nd smallest & 1.539 & 0.601 \\
3rd smallest & 77.2 & 21.7 \\
\bottomrule
\end{tabular}
\end{center}

The separations vary smoothly and monotonically with~$\delta$.
No discontinuity, cusp, or topological transition occurs at
$\delta = 2/9$.


\section{Democratic limit: exact factorisation}
\label{app:democratic}

For $m_1 = m_2 = m_3 = m$, the Seiberg--Witten sextic
factors as~\eqref{eq:dem_factor}.
The bracket cubic
$g(t) = t^3 - 6mt^2 + 9m^2t - 4\Lambda^3$
has roots:

\begin{itemize}
  \item For $\Lambda < m$: three distinct real roots.
    The curve is smooth.
  \item At $\Lambda = m$:
    $g(t) = (t - m)^2(t - 4m)$.
    The double root at $t = m$ ($x = 0$) signals a
    \textbf{massless BPS monopole}.
  \item For $\Lambda > m$: one real root and two complex conjugate
    roots.
\end{itemize}

The full root structure at $\Lambda = m$:
$x = -m$ (multiplicity~3), $x = 0$ (multiplicity~2),
$x = 3m$ (multiplicity~1).
The genus-2 curve degenerates to genus~1 (one cycle pinches).

In the $\mathcal{N}{=}1$ theory obtained by giving the adjoint
a mass~$\mu_\Phi$, the singularity at $\Lambda = m$ is
resolved: the monopole acquires a mass of order~$\mu_\Phi$ and
the vacuum is shifted away from the singular point.
The $\mathcal{N}{=}2 \to \mathcal{N}{=}1$ deformation thus
provides a dynamical mechanism for lifting the monopole singularity.


\section{K\"ahler protection estimate}
\label{app:kahler}

In the $\mathcal{N}{=}2$ limit, the K\"ahler metric on the Higgs
branch is one-loop exact.
The leading deviation from generation-universality after
$\mathcal{N}{=}2 \to \mathcal{N}{=}1$ breaking comes from the
adjoint-mass deformation and is proportional to the mass splitting:
\begin{equation}
  \frac{\delta Z_i}{Z}
  \sim \frac{g^2}{16\pi^2}\,
  \frac{m_i^2 - \bar{m}^2}{\mu_\Phi^2}\,
  \ln\frac{\mu_\Phi}{\Lambda}\,,
\end{equation}
where $\bar{m}^2 = \frac{1}{3}\sum_j m_j^2$ and the loop factor
arises from the gaugino--adjoint propagator.

For the lepton sector:
\begin{align}
  \frac{m_\tau^2 - \bar{m}^2}{\mu_\Phi^2}
  &\approx \frac{m_\tau^2}{\mu_\Phi^2}\,,\\
  \frac{m_e^2 - \bar{m}^2}{\mu_\Phi^2}
  &\approx -\frac{m_\tau^2}{3\mu_\Phi^2}\,.
\end{align}
The ratio $\delta Z_\tau / \delta Z_e \approx -3$, so the leading
correction to the Koide ratio is
\begin{equation}
  \delta K
  \approx \frac{\delta Z_\tau - \delta Z_e}{Z}
  \sim \frac{g^2}{4\pi^2}\,
  \frac{m_\tau^2}{\mu_\Phi^2}\,
  \ln\frac{\mu_\Phi}{\Lambda}\,.
  \label{eq:deltaK_estimate}
\end{equation}
The value $\mu_\Phi \approx 4.5\TeV$ is not a free parameter:
it is the critical scale below which the quartic~\eqref{eq:quartic}
is strong enough to sustain the mixed vacuum, determined by the
RG crossing condition $K^{-1}(d,s,b) = 2/3$.
For $g^2/(4\pi) \sim 0.3$, $m_\tau \approx 1.78\GeV$,
$\mu_\Phi \sim 4.5\TeV$, and $\ln(\mu_\Phi/\Lambda) \sim 10$:
\begin{equation}
  \delta K
  \sim \frac{0.3}{4\pi}\times
  \frac{(1.78)^2}{(4500)^2}\times 10
  \sim 10^{-9}\,,
\end{equation}
which is well below the observed
$\delta K_{\rm exp} \approx 5 \times 10^{-5}$.
The discrepancy suggests that the dominant contribution to
$\delta K$ comes from a source other than the adjoint-mass
deformation---most likely from the running of quark masses
between the family confinement scale and the pole mass
scale~\cite{Sumino:2009,MartinRobertson:2019}.


\section{$\UU(3)$ charge algebra and uniqueness of the quadratic power}
\label{app:charges}

Let $\UU(3)$ be a family symmetry with generator
$Q_F = z_0\,I + T$, where $T = \diag(z_1, z_2, z_3)$ lies in the
Cartan of $\SU(3)$ with $\sum_i z_i = 0$.
Canonical normalisation fixes
\begin{equation}
  z_0 = \frac{1}{\sqrt{6}}\,,
  \qquad
  \sum_{i=1}^{3} z_i^2 = \frac{1}{2}\,,
  \label{eq:canon}
\end{equation}
so that the eigenvalues $\zch_i \equiv z_0 + z_i$ satisfy
$\sum_i \zch_i = \sqrt{3/2}$ and $\sum_i \zch_i^2 = 1$.
An explicit parametrisation of the Cartan direction is
\begin{equation}
  z_i = \frac{1}{\sqrt{3}}
  \cos\!\left(\delta + \frac{2\pi(i-1)}{3}\right),
  \quad i = 1,2,3\,.
\end{equation}

\paragraph{Charge-squared mass law.}
If $m_i = \mu\,\zch_i^2$, the Koide ratio becomes
\begin{equation}
  K = \frac{\sum_i \zch_i^2}{(\sum_i |\zch_i|)^2}
    = \frac{1}{3/2} = \frac{2}{3}\,,
\end{equation}
independent of~$\delta$.

\begin{theorem}[Uniqueness of the quadratic]
\label{thm:unique}
For integer $n$, define
$K_n(\delta) = \sum_i \zch_i^n / (\sum_i \zch_i^{n/2})^2$
whenever the ratio is real.
Then $K_n$ is independent of the Cartan direction~$\delta$
if and only if $n = 2$ (besides the trivial $n = 0$).
\end{theorem}

\begin{proof}
The binomial expansion of $\sum_i (z_0 + z_i)^n$ contains the
cubic Casimir $\sum_i z_i^3 = (1/\sqrt{3})^3\cos 3\delta/(3/4)
\propto \cos 3\delta$ for $n \geq 3$.
For $n = 2$: $\sum_i \zch_i^2 = 3z_0^2 + \sum z_i^2 = 1$
(no $\delta$-dependence).
For $n = 1$: $\sum_i \zch_i = 3z_0 = \sqrt{3/2}$
(trivially constant, but $K_1 = 1/3$ requires $\zch_i \geq 0$
and is not the Koide ratio).
For $n \geq 3$: the $\cos 3\delta$ term from the cubic Casimir
makes $\sum_i \zch_i^n$ depend on~$\delta$, so $K_n$ varies.
\end{proof}


\section{The interleaving theorem}
\label{app:interleaving}

The two basic hierarchies produced by the Seiberg seesaw are the
direct and inverse branches:
\begin{equation}
  a_i = \mu\,\zch_i^2\,,
  \qquad
  d_i = \nu\,\zch_i^{-2}\,,
\end{equation}
with opposite orderings for $\zch_1 < \zch_2 < \zch_3$:
$a_1 < a_2 < a_3$ while $d_1 > d_2 > d_3$.

\begin{theorem}[Interleaving]
\label{thm:interleaving}
The six unmixed masses $\{a_1, a_2, a_3, d_1, d_2, d_3\}$
form the alternating ladder
\begin{equation}
  d_1 > a_3 > d_2 > a_2 > d_3 > a_1
\end{equation}
if and only if
\begin{equation}
  \max(\zch_1^2 \zch_3^2,\, \zch_2^4)
  < \frac{\nu}{\mu}
  < \zch_2^2\,\zch_3^2\,.
  \label{eq:window}
\end{equation}
\end{theorem}

\begin{proof}
The adjacent inequalities reduce to three independent conditions:
$\nu/\mu > \zch_1^2\zch_3^2$,
$\nu/\mu < \zch_2^2\zch_3^2$, and
$\nu/\mu > \zch_2^4$.
The remaining inequalities follow from
$\zch_1 < \zch_2 < \zch_3$.
\end{proof}

Each consecutive triple in the ladder mixes direct and inverse
entries, producing the quark mass ``waterfall'' observed
empirically~\cite{Zenczykowski:2012}.
The quartic correction~\eqref{eq:quartic} perturbs the eigenvalues
away from the unmixed values $a_i, d_i$, and its hierarchy of
corrections (largest for the lightest family, smallest for the
heaviest) translates into a hierarchy of deviations from $K = 2/3$
in mixed waterfall windows.


%======================================================================
\begin{thebibliography}{99}

\bibitem{Chew:1962}
  G.F.~Chew,
  Rev.\ Mod.\ Phys.\ \textbf{34}, 394 (1962).

\bibitem{Rivero:2024epjc}
  A.~Rivero,
  Eur.\ Phys.\ J.\ C \textbf{84}, 1058 (2024),
  arXiv:2407.05397.

\bibitem{Kaplan:1991dc}
  D.B.~Kaplan,
  Nucl.\ Phys.\ B \textbf{365}, 259 (1991).

\bibitem{HananyWitten:1997}
  A.~Hanany and E.~Witten,
  Nucl.\ Phys.\ B \textbf{492}, 152 (1997),
  arXiv:hep-th/9611230.

\bibitem{Witten:1997sc}
  E.~Witten,
  Nucl.\ Phys.\ B \textbf{500}, 3 (1997),
  arXiv:hep-th/9703166.

\bibitem{Seiberg:1994bz}
  N.~Seiberg,
  Nucl.\ Phys.\ B \textbf{435}, 129 (1995),
  arXiv:hep-ph/9411149.

\bibitem{Csaki:1996zb}
  C.~Cs\'aki, M.~Schmaltz, and W.~Skiba,
  Phys.\ Rev.\ Lett.\ \textbf{78}, 799 (1997),
  arXiv:hep-th/9610139.

\bibitem{IntriligatorSeibergShenker:1995}
  K.~Intriligator, N.~Seiberg, and S.H.~Shenker,
  Phys.\ Lett.\ B \textbf{342}, 152 (1995),
  arXiv:hep-ph/9410203.

\bibitem{SeibergWitten:1994a}
  N.~Seiberg and E.~Witten,
  Nucl.\ Phys.\ B \textbf{426}, 19 (1994),
  arXiv:hep-th/9407087.

\bibitem{Fayet:1976}
  P.~Fayet,
  Nucl.\ Phys.\ B \textbf{113}, 135 (1976).

\bibitem{HMS:1997}
  T.~Hirayama, N.~Maekawa, and S.~Sugimoto,
  arXiv:hep-th/9705069 (1997).

\bibitem{ArgyresDouglas:1995}
  P.C.~Argyres and M.R.~Douglas,
  Nucl.\ Phys.\ B \textbf{448}, 93 (1995),
  arXiv:hep-th/9505062.

\bibitem{Kostant:1959}
  B.~Kostant,
  Amer.\ J.\ Math.\ \textbf{81}, 973 (1959).

\bibitem{SeibergWitten:1994b}
  N.~Seiberg and E.~Witten,
  Nucl.\ Phys.\ B \textbf{431}, 484 (1994),
  arXiv:hep-th/9408099.

\bibitem{ArgyresPlesserShapere:1996}
  P.C.~Argyres, M.R.~Plesser, and A.D.~Shapere,
  Phys.\ Rev.\ Lett.\ \textbf{75}, 1699 (1995),
  arXiv:hep-th/9505100.

\bibitem{Townsend:1995}
  P.K.~Townsend,
  Phys.\ Lett.\ B \textbf{350}, 184 (1995),
  arXiv:hep-th/9501068.

\bibitem{Schwarz:1995}
  J.H.~Schwarz,
  Phys.\ Lett.\ B \textbf{360}, 13 (1995),
  arXiv:hep-th/9508143.

\bibitem{Strominger:1995}
  A.~Strominger,
  Phys.\ Lett.\ B \textbf{383}, 44 (1996),
  arXiv:hep-th/9512059.

\bibitem{Klemm:1996}
  A.~Klemm, W.~Lerche, S.~Yankielowicz, and S.~Theisen,
  Phys.\ Lett.\ B \textbf{344}, 169 (1995),
  arXiv:hep-th/9411048.

\bibitem{ArgyresShapere:1995}
  P.C.~Argyres and A.E.~Faraggi,
  Phys.\ Rev.\ Lett.\ \textbf{74}, 3931 (1995),
  arXiv:hep-th/9411057.

\bibitem{Zenczykowski:2012}
  P.~Zenczykowski,
  Phys.\ Rev.\ D \textbf{86}, 117303 (2012),
  arXiv:1210.4125.

\bibitem{Sumino:2009}
  Y.~Sumino,
  JHEP \textbf{05}, 075 (2009),
  arXiv:0812.2103.

\bibitem{MartinRobertson:2019}
  S.P.~Martin and D.G.~Robertson,
  Phys.\ Rev.\ D \textbf{100}, 073004 (2019),
  arXiv:1907.02500.

\bibitem{Koide:1983qe}
  Y.~Koide,
  Phys.\ Lett.\ B \textbf{120}, 161 (1983).

\bibitem{MasieroVeneziano:1985}
  A.~Masiero and G.~Veneziano,
  Nucl.\ Phys.\ B \textbf{249}, 593 (1985).

\bibitem{Antusch:2025ahs}
  S.~Antusch, K.~Hinze, and S.~Saad,
  arXiv:2510.01312 (2025).

\end{thebibliography}

\end{document}
